\documentclass[lualatex,paper=b5j,fontsize=10pt,twoside]{jlreq}

\usepackage{surikumi}

\MathMagazineSetup{
  feature={},
  series={要点の整理／数学 B},
  series-position=left,
  pattern=scales,
  title={数列の確認},
  author={数学編集部},
  author-font=mincho,
  author-position=right,
  title-fill=black,
  deck={演習前の確認用に，基本公式と考え方を短く整理します．}
}

\begin{document}

\MakeMathMagazineTitle
\MMConfirmationReset

\MMConfirmationSection{基本的な数列}
\MMConfirmationSubsection{等差数列}

初項 $a$，公差 $d$ の等差数列 $\{a_n\}$ について，次の公式を確認する．

\begin{mmconfirmationroman}
  \item $a_{n+1}-a_n=d$
  \item $a_n=a+(n-1)d$
  \item 和 $S_n=\sum_{k=1}^{n}a_k$ は
    \[
      S_n=\dfrac{a+a_n}{2}n
    \]
\end{mmconfirmationroman}

\MMConfirmationGeneralComment{一般項では，初項から第 $n$ 項までに公差を $n-1$ 回加える．}

\MMConfirmationSubsection{等比数列}

初項 $a$，公比 $r$ の等比数列 $\{a_n\}$ について，

\begin{mmconfirmationroman}
  \item $a_{n+1}=ra_n$
  \item $a_n=ar^{n-1}$
  \item $r\neq 1$ のとき
    \[
      S_n=\dfrac{a(r^n-1)}{r-1}
    \]
\end{mmconfirmationroman}

\MMConfirmationGeneralNote{$r=1$ のときは，$S_n=na$ と別に扱う．}

\MMConfirmationAdvancedComment{和の変形は，$\sum_{k=1}^{n}kr^k$ の計算にも応用できる．}

\MMConfirmationAdvancedNote{等比数列の和を $r$ で微分する方法もある．}

\MMConfirmationSection{数列の和}
\MMConfirmationSubsection{和から一般項へ}

数列 $\{a_n\}$ の初項から第 $n$ 項までの和を $S_n$ とする．

\MMConfirmationRunInHeading{基本公式}
$a_1=S_1$ であり，$n\geq 2$ のとき
\[
  a_n=S_n-S_{n-1}
\]
である．

\MMConfirmationGeneralNote[1]{$n\geq 2$ では，$S_n-S_{n-1}$ を計算する．}
\MMConfirmationGeneralNote[2]{$n=1$ は，$a_1=S_1$ に戻って確認する．}

\MMConfirmationSubsection{階差と和}

数列 $\{b_n\}$ が
\[
  b_{n+1}-b_n=a_n
\]
を満たすとする．$n\geq 2$ のとき，
\[
  b_n=b_1+\sum_{k=1}^{n-1}a_k
\]
となる．

\begin{mmconfirmationproof}
等式 $b_{k+1}-b_k=a_k$ を $k=1,2,\ldots,n-1$ について加えると，
\begin{align*}
  \sum_{k=1}^{n-1}(b_{k+1}-b_k)
    &= (b_2-b_1)+\cdots+(b_n-b_{n-1}) \\
    &= b_n-b_1
\end{align*}
となる．
\end{mmconfirmationproof}

\end{document}
